\section{Resumen}

En este informe se realizaron mediciones sobre circuitos de corriente alterna monofásica utilizando distintos instrumentos, entre ellos el vatímetro, amperímetro y voltímetro. A partir de las mediciones efectuadas se analizaron las relaciones entre potencia activa, reactiva y aparente, así como el comportamiento del factor de potencia en distintas configuraciones del circuito.

Durante la experiencia se trabajó con una carga resistiva formada por lámparas incandescentes, una bobina y un banco de capacitores conectable en paralelo. Asimismo, se estudió el efecto producido por la incorporación de un núcleo ferromagnético en la inductancia y la influencia de los capacitores sobre la compensación de potencia reactiva y la corrección del factor de potencia.

En general, los resultados experimentales obtenidos presentaron una buena concordancia con los valores teóricos esperados, permitiendo verificar el comportamiento de los circuitos en corriente alterna, analizar la compensación del factor de potencia y estudiar la influencia de distintos núcleos sobre las características eléctricas de la bobina.

\newpage
